\documentclass[10pt,a4paper]{extarticle}
\pagenumbering{gobble}
\usepackage[left=0.6in,top=0.6in,right=0.6in,bottom=0.6in]{geometry}
\usepackage{hyperref}
\usepackage{titlesec}
\usepackage{enumitem}
\usepackage{graphicx}
\titleformat{\section}{\large\bfseries\scshape\raggedright}{}{0em}{}[\titlerule]
\titlespacing*{\section}{0pt}{12pt}{8pt}
\begin{document}
\begin{center}
    \begin{minipage}{\textwidth}
        \centering
        {\LARGE\textbf{Umut Albayrak}} \hspace{2pt} {\large{Backend Developer | Bilgisayar Mühendisliği Öğrencisi}}\\[10pt]
        \href{mailto:umutalbayrakk24@gmail.com}{umutalbayrakk24@gmail.com} \textbullet\
        \href{https://umutalbayrak.dev}{umutalbayrak.dev} \textbullet\
        \href{https://linkedin.com/in/umut-albayrak24}{linkedin.com/in/umut-albayrak24} \textbullet\
        \href{https://github.com/Umutalb}{github.com/Umutalb}
    \end{minipage}
\end{center}
\section{Özet}
C\# ve ASP.NET Core ile backend geliştirmeye odaklanan Bilgisayar Mühendisliği öğrencisi. RESTful API geliştirme, veritabanı tasarımı ve clean architecture prensiplerini uygulama konusunda yetkin. Git, Firestore ve çevik (agile) ekip çalışması deneyimine sahip. Performans, ölçeklenebilirlik ve güvenliğe odaklanan Junior Backend Developer olarak katkı sağlamak istiyor.
\section{Teknik Deneyim}
\textit{DatingApp2025} \hspace{3cm} \textit{Backend Developer \& Scrum Master}\\
\begin{itemize}[leftmargin=*,noitemsep,topsep=0pt]
    \item 5 kişilik ekip projesinde hem Backend Developer hem de Scrum Master rolünü üstlendi.
    \item ASP.NET Core Web API ile kimlik doğrulama, raporlama, eşleştirme ve bildirim modüllerini geliştirdi.
    \item Google Firestore entegrasyonu ile verilerin kalıcı şekilde saklanmasını sağladı.
    \item Scrum Master olarak sprint planlama ve günlük toplantıları yönetti, ekip içi iş birliğini ve görev takibini sağladı.
    \item Frontend geliştirici (React) ve AI ekibi (Flask) ile iş birliği yaparak öneri sisteminin entegrasyonunu ve API kullanımını kolaylaştırdı.
    \item API dokümantasyonu ve GitHub üzerinden versiyon kontrolüne katkıda bulundu.
\end{itemize}
\vspace{0.5cm} 
\textit{PhoneBook} \hspace{3cm} \textit{C++ Konsol Uygulaması}
\begin{itemize}[leftmargin=*,noitemsep,topsep=0pt]
    \item CRUD işlemleri ve File I/O ile C++ ve STL kullanılarak geliştirildi.
    \item Veriler, yükleme/kaydetme güvenilirliği için \texttt{fstream} ile CSV-benzeri formatta saklandı.
\end{itemize}
\vspace{0.5cm} 
\textit{JobApplicationChecker} \hspace{3cm} \textit{C++ Konsol Uygulaması}\\
\begin{itemize}[leftmargin=*,noitemsep,topsep=0pt]
    \item Yaş, ehliyet ve üniversite mezuniyeti kriterlerine göre kural tabanlı uygunluk kontrolü yaptı.
    \item İç içe koşullar ve büyük/küçük harf duyarsız string kontrolleriyle net bir kontrol akışı uygulandı.
    \item Kullanıcı yönlendirmeleri ve hata mesajlarına odaklanarak basit bir işe alım senaryosu tasarlandı.
\end{itemize}
\section{Beceriler}
\textbf{Programlama Dilleri:} C\#, C++ \\
\textbf{Backend Frameworks \& Runtimes:} ASP.NET Core, Entity Framework Core, RESTful API Development \\
\textbf{Frontend:} HTML, CSS \\
\textbf{Veritabanları:} MSSQL Server, Google Firestore \\
\textbf{Diğer:} Git, GitHub, REST API, Postman, Swagger \\
\textbf{Yöntemler:} Scrum, Agile, Clean Architecture
\section{Eğitim}
\textbf{Van Yüzüncü Yıl Üniversitesi} \hfill Van, Türkiye\\
\textit{Bilgisayar Mühendisliği Lisans Programı} \hfill Eki 2023 -- Haz 2027\\
\section{Sertifikalar}
\begin{itemize}[leftmargin=*,noitemsep,topsep=0pt]
    \item HackerRank C\# (Basic) Sertifikası – 2025
\end{itemize}
\section{Diller}
\begin{itemize}[leftmargin=*,noitemsep,topsep=0pt]
    \item Türkçe – Ana Dil
    \item İngilizce – Başlangıç Seviyesi
\end{itemize}
\end{document}
